\documentclass[11pt]{article}
\usepackage[margin=1in]{geometry}
\usepackage{amsmath,amssymb}
\usepackage[T1]{fontenc}
\usepackage{hyperref}

\title{Quaternary Quartic Rank-7 Exploration Log}
\author{}
\date{}

\begin{document}
\maketitle

\section*{2026-05-04: Rank-deficient branch formalized}

\paragraph{Starting state.}
The Lean root theorem
\texttt{QuaternaryQuartic.quaternaryQuartic\_rankSeven\_no\_spurious\_socp}
was still proved through the scaffold axiom
\texttt{quaternaryQuartic\_rankSeven\_no\_spurious\_socp\_placeholder}.
The full verification harness therefore built the files but failed at
\texttt{\#print axioms}, as expected.

\paragraph{Largest gap selected.}
The main blueprint splits first on whether the seven quadratic coordinates
\(u_i\) are linearly dependent.  The rank-deficient branch is independent of
the low-rank catalecticant support theorem and can be closed directly from the
SOCP inequalities and the SOS certificate.  This made it the largest remaining
gap that could be reduced without first formalizing the apolar Hilbert-function
argument.

\paragraph{Lean progress.}
Added \texttt{QuaternaryQuarticProof/Certificate.lean}.  It defines:
\[
\texttt{relationDirection}\ c\ q=(i\mapsto c_i q),\qquad
\texttt{relationPoly}\ u\ c=\sum_i c_i u_i,
\]
plus the supporting admissibility, \(A\)-image, and \(\sigma\)-formula lemmas.
The main new theorem is
\[
\texttt{residual\_eq\_zero\_of\_constant\_relation}.
\]
It proves that if there is a nonzero scalar relation
\(\sum_i c_i u_i=0\), then any SOS/SOCP point has
\(\texttt{residual}\ p\ u=0\).

\paragraph{Proof idea.}
For each SOS summand \(q_j^2\), use the kernel direction
\(w_i=c_i q_j\).  Then \(A_u(w)=0\) and
\[
\sigma(w)=\left(\sum_i c_i^2\right)q_j^2.
\]
The Hessian inequality gives
\(\left(\sum_i c_i^2\right)B(q_j^2,R)\ge0\), and the scalar is positive
because \(c\ne0\).  Summing over the SOS certificate gives \(B(p,R)\ge0\).
FOCP at \(v=u\) gives \(B(\sigma(u),R)=0\), hence the objective equals
\(-B(p,R)\le0\).  Positive-definiteness gives the opposite inequality, so the
objective is zero and the residual vanishes.

\paragraph{Verification.}
The new module and root proof file typecheck:
\[
\texttt{lake env lean QuaternaryQuarticProof/Certificate.lean},\qquad
\texttt{lake env lean QuaternaryQuarticProof.lean}.
\]
The full harness command available in the repository is
\texttt{./scripts/verify\_quaternary\_quartic.sh}; the user-provided
\texttt{verify\_quarternary\_quartic.sh} path does not exist.  The full harness
still fails only because the root theorem depends on the existing placeholder
axiom.

\paragraph{Remaining blockers.}
The largest remaining gap is the full-rank branch.  The next proof-serving
steps are:
\begin{enumerate}
\item formalize the catalecticant kernel \(K_\lambda\) and prove
  \(L_u\subseteq K_\lambda\) from FOCP;
\item express full rank of the seven coordinates as
  \(\ker(\texttt{relationPolyLin}\ u)=\bot\), yielding \(\dim L_u=7\);
\item formalize the low-rank catalecticant support theorem and the rank
  \(1,2,3\) syzygy-to-negative-curvature constructions.
\end{enumerate}

No Julia numerical claims were used in this round.

\end{document}
